\section{Digit factorials, normal tests, and admissible assembly}
\label{sec:digit-tests}

\status{We prove the factorial-unit approximation, construct the normal test, and perform the coefficient selection with its exact scalar. We then prove outgoing symbol membership under the stated frame, degree and margin data. The incoming assertion \texttt{MemAt}, placing the whole operator-word family in the signed-chain space, remains a separate input. The expanded DT note gives the segment-by-segment membership calculation.}

\subsection{A shorter factorial-unit approximation}

Let $p$ be any prime, $P_v=p^{g_v}$, $0\le\rho_v<P_v$, and $\mathbf b\in\mathbb N^V$, for a finite set $V$. Put
\[
 \Gamma(\mathbf b)=\prod_{v\in V}
 \frac{(\rho_v+P_vb_v)!}{\rho_v!\,b_v!\,(P_v!)^{b_v}}.
 \tag{DT.1}
\]
This is a $p$-adic unit. For every integer $c\ge1$ and $\varepsilon\in\{1,-1\}$, there is a polynomial
\[
 Q(\mathbf b)=\sum_{|\mathbf j|\le2(c-1)}q_{\mathbf j}
       \prod_v\binom{b_v}{j_v},\qquad q_{\mathbf j}\in\mathbb Z_{(p)},
 \qquad v_p(\Gamma(\mathbf b)^\varepsilon-Q(\mathbf b))\ge c.
 \tag{DT.2}
\]
The bound is independent of $|V|$ and holds also at $p=2,3$. It improves the source interfaces $4c$ and $6c$; later numerical constants are left unchanged.

Here is a finite-product proof. If $t_i\ge1$ and $T=\sum_i t_i$, the ordered sum
\[
 S_{\mathbf t}(b)=\sum_{0\le u_1<\cdots<u_m<b}u_1^{t_1}\cdots u_m^{t_m}
\]
is an integer-valued polynomial of degree at most $T+m\le2T$. Induct using discrete summation and $\sum_{u<b}\binom uj=\binom b{j+1}$; Newton coefficients are integer finite differences. Thus a factor $F(u)=1+\sum_{t\ge1}a_tu^t$ with $a_t\in p^t\mathbb Z_{(p)}$ has
$\prod_{u<b}F(u)$ approximated modulo $p^c$ by keeping total grade $T<c$. The contribution of grade $T$ is an integral binomial polynomial of degree at most $2T$, with coefficients divisible by $p^T$. Products over several variables retain this property by adding grades. \textsf{[DT1--DT2]}

Write $U(n)=\prod_{1\le i\le n,\ p\nmid i}i$. Removing multiples of $p$ repeatedly gives
\[
 \frac{(\rho+Pb)!}{\rho!b!(P!)^b}
 =\prod_{i=1}^{\rho}\left(1+\frac{P}{i}b\right)
  \prod_{\ell=1}^{g}\prod_{u=0}^{b-1}
       \prod_{\substack{1\le a\le p^\ell\\p\nmid a}}
            \left(1+\frac{p^\ell}{a}u\right).
 \tag{DT.3}
\]
Every nonconstant linear coefficient is in $p\mathbb Z_{(p)}$. Apply the graded product argument; for the inverse use the finite geometric expansion modulo $p^c$. This proves (DT.2), including uniformity in all natural $\mathbf b$, without an infinite-product convergence argument. \textsf{[DT3--DT4]}

\subsection{Normal projection with its exact scalar}

The coefficient multiplier $\prod_v\binom{\alpha_v}{j_v}$ is the operator
$X^{\mathbf j}\partial^{[\mathbf j]}$. Its conjugation module has all coordinate weights at most $|\mathbf j|$, so it increases a prefix's largest highest-weight coordinate by at most that amount. It also preserves each letter degree. This proves the weight cost of (DT.2), without an integral isotypic decomposition. \textsf{[DT5--DT6]}

In the actual local setting take $w^2=p$. On the normal polynomial space let $e$ be the joint projector onto the chosen prefix types, $D$ its summed IW4 denominator, and $\nu$ a monomial of normal multidegree $d^{\rm no}$. Define
\[
 F_{\rm no}=v_p\!\left(\prod_xd_x^{\rm no}!\right),\quad U=v_p(\nu!),\quad
 c_Y=w^{F_{\rm no}}p^Dw^{-U},\qquad Y=c_Ye(X^\nu).
 \tag{DT.4}
\]
The polynomial $w^{-U}X^\nu$ belongs to the weighted Fock lattice. Projector integrality followed by the factor $w^{F_{\rm no}}$ makes $Y$ coefficient-integral. Apply (DT.2) to the inverse normal factorial unit and set $N_0=\operatorname{diag}(Q)Y$. Then
\[
 N_0-\operatorname{diag}(\Gamma_{\rm no}^{-1})Y
       \in p^c\mathcal O[X],\qquad
 v(c_Y)=F_{\rm no}/2+D-U/2.
 \tag{DT.5}
\]
Its prefix first weights increase by at most $2(c-1)\le4c$, exactly the allowed \texttt{NormTest} slack. Thus the substantive \texttt{NoAt} construction is explicit. \textsf{[DT7]}

\subsection{Selecting the plane coefficient and preserving detection}

Start with a coefficient-integral polynomial $G$ having a unit coefficient at $\alpha_0$. Split its true exponents uniquely as $\rho+P\beta$, and separate plane and normal carries. Put
$V_P=v_p(\beta_{P,0}!)$ and $V_N=v_p(\beta_{N,0}!)$.
For the resulting normal coefficient polynomial $y$, $b(y)\le V_N/2$.
Since the finite joint projectors sum to the identity, one satisfies $b(ey)\le b(y)$.
Choose a normal monomial $\nu$ attaining this norm, and let $P_c$ be its plane coefficient after projection. Then
\[
 b(P_c)+v_p(\nu!)/2\le(V_P+V_N)/2.
\]
Choose $z$ of valuation $b(P_c)$ from an attaining coefficient and the appropriate power of $w$; set $F_2=z^{-1}P_c$, so $b(F_2)=0$.
The incoming \texttt{MemAt} condition is precisely what puts $F_2$ in the required signed-chain space. \textsf{[DT8]}

The exact factorial identity on this digit block is
\[
 \alpha!=C_\rho\,\beta_P!\,\beta_N!\,
       \Gamma_{\rm pl}(\beta_P)\Gamma_{\rm no}(\beta_N),\qquad
 C_\rho=\prod_v\rho_v!\prod_x(P_x!)^{d_x^{\rm pl}+d_x^{\rm no}}.
\]
Fock self-adjointness therefore gives the ideal pairing scalar
$k_0=C_\rho c_Y\nu!z$. With $F_{\rm tot}=v_p(\prod_xd_x!)$ and
$F_{\rm pl}=v_p(\prod_xd_x^{\rm pl}!)$, the explicit size condition
\[
 D+(F_{\rm pl}-V_P)/2+(F_{\rm no}-V_N)/2\le D_n
\]
implies
\[
 v(k_0)\le F_{\rm tot}-F_{\rm pl}/2+D_n.
 \tag{DT.6}
\]
Indeed $v(C_\rho)+V_P+V_N=v_p(\alpha_0!)\le F_{\rm tot}$. Replacing the ideal normal factor by $N_0$ introduces pairing error of valuation at least $c+F_{\rm tot}-\kappa_{\rm fl}$. These are the exact \texttt{PickAt} estimates, and CL8 already gives $D_n\le9Nk^2\log_p a$.

\subsection{Why the output is an allowed symbol}

This last step concerns \texttt{MembAt}, distinct from \texttt{MemAt}. Under the source's \texttt{MFH} frame and degree conditions, each relevant slack satisfies $s\ge128c$.
At a cut put $T=\texttt{thStar}$ and $\tau=T-s/2-60c$.
PG's chain and link bounds give plane order at least $\tau$; the normal first weight is at most $a+4c$.
For an even position the definition of $T$ gives $T\ge a_{\rm after}+m+s$, hence $\tau\ge a_{\rm after}+m+4c$. At an odd position use the next cut with $a_{\rm before}$. Also $T\ge B/2$ gives $\tau\ge B/2-s$. These are exactly the two separation bounds defining the segment symbol space. \textsf{[DT10]}

All segment constraints cut out rational subspaces. Their intersections with the coefficient lattices are saturated and therefore direct summands over the DVR. Extending their bases proves that simultaneous constraints on disjoint segments give an integral sum of products of constrained segment factors. No integral decomposition into irreducible representations is needed. \textsf{[DT9]}
The carry identity $p^{z_j-x}p^{z_{j+1}-z_j}=p^{z_{j+1}-x}$ verifies the actual Frobenius substitution on variables. Goodness and \texttt{BlockOK} give the physical degrees; $\rho<P$ gives the label bounds. The normal order is fixed, with plane/normal roles exchanged at odd positions. Thus the assembled test is in the permitted integral symbol span. The DT report proves all four interfaces \texttt{DecAt}, \texttt{AlgAt}, \texttt{DegAt}, and \texttt{LabAt} pointwise. \textsf{[DT10]}

The initial-symbol and operator-word proof of incoming \texttt{MemAt}, the quantitative selected-symbol dimension, and the full numerical frame supply remain to be reconstructed. Closing the outgoing membership calculation does not close these different assertions. \textsf{[DT11]}
